\subsection{从能谱相干算出电子脉冲列}
\label{sec:electron-temporal-bunching-worked}

在能谱上看到一串边带以后，还不能仅凭各峰的高度确定电子的时间分布；不同边带的相位决定它们在何时相长干涉。下面把上一节的密度矩阵传播式用于 PINEM 输出，求出周期密度及其 Fourier 谐波。令 $\theta_\zeta=\kappa_z\zeta$ 表示一个空间周期内的位置，$s_{\rm d}=\omega_{\rm r}t\simeq\omega_{\rm r}L/v_0$ 表示传播中积累的反冲相位；两者均无量纲。相对周期密度为
\begin{equation}
 \mathfrak n_e(\theta_\zeta,s_{\rm d})
 =\sum_{\ell m}\rho_{\ell m}^{(+)}
 \ee^{\ii(\ell-m)\theta_\zeta-\ii(\ell^2-m^2)s_{\rm d}},\qquad
 \frac1{2\pi}\int_0^{2\pi}\mathfrak n_e\dd\theta_\zeta=1.
 \label{eq:electron-bunching-normalized-density}
\end{equation}
$\mathfrak n_e$ 是相对于慢包络背景的无量纲密度，不是单位体积内的电子数。在固定探测面，相对于中心电子的到达延迟为 $\tau_{\rm d}=t-L/v_0$，所以 $\theta_\zeta=-\omega\tau_{\rm d}$；正的到达延迟对应负的共动坐标。把 $s_{\rm d}$ 在一个光学周期内视为常数，还要求实际占据阶数上 $\ell^2\omega_{\rm r}|\tau_{\rm d}|\ll1$。

对任意混态，$\rho_{\ell m}^{(+)}$ 的对角部分给常背景，非对角部分才给随 $\theta_\zeta$ 变化的项。若能谱为 $P_\ell$ 而相位完全随机，则 $\rho_{\ell m}=P_\ell\delta_{\ell m}$，无论传播多远都有 $\mathfrak n_e=1$。相反，纯 PINEM 态的 $c_\ell=J_\ell(2|\beta|)\ee^{\ii\ell\varphi_\beta}$ 保留确定的交叉项。刚离开相互作用区时
\begin{equation*}
 \sum_\ell c_\ell\ee^{\ii\ell\theta_\zeta}
 =\exp[2\ii|\beta|\sin(\theta_\zeta+\varphi_\beta)],
\end{equation*}
模平方仍为一；后续二次色散才把相位调制变为密度调制。产生能谱边带与产生时间脉冲是不同步骤。

将周期密度写为 $\mathfrak n_e=\sum_d B_d\ee^{\ii d\theta_\zeta}$，其 Fourier 系数为
\[
B_d=\frac1{2\pi}\int_0^{2\pi}\dd\theta_\zeta\,
\mathfrak n_e(\theta_\zeta,s_{\rm d})\ee^{-\ii d\theta_\zeta}.
\]
把密度的双重求和代入，角积分只留下 $\ell-m=d$ 的项。这解释了为什么第 $d$ 阶密度谐波需要所有相隔 $d$ 条边带的相干共同贡献。对上述纯 PINEM 输入，Bessel 和式进一步给出
\begin{equation}
 B_d(s_{\rm d})=(-\ii)^d\ee^{\ii d\varphi_\beta}
 J_d[4|\beta|\sin(ds_{\rm d})],\qquad d\in\mathbb Z.
 \label{eq:electron-pinem-bunching-closed}
\end{equation}
$d$ 是被读取的密度谐波阶数，不是某一个末态边带编号；$B_d$ 汇总所有相隔 $d$ 阶的相干。该式保持 $B_0=1$、$B_{-d}=B_d^*$，并在 $s_{\rm d}=0$ 时使所有非零阶谐波消失。

\begin{exbox}[第一谐波最强的距离：弱调制与强调制]
给定 $b=|\beta|>0$，求第一谐波首次达到最大值的漂移距离。由闭式，问题化为在 $0\le s_{\rm d}\le\pi/2$ 内最大化
$|B_1|=|J_1(4b\sin s_{\rm d})|$。

令 $x_*$ 为 $J_1'(x)=0$ 的第一个正根。递推关系给
$J_1'(x)=[J_0(x)-J_2(x)]/2$，所以该常数也由 $J_0(x_*)=J_2(x_*)$ 定义；$x_*\simeq1.84118$。$J_1$ 的首个正极大值也是其绝对最大值。现在只需检查驱动能否达到这个参数：
\begin{equation*}
 s_{\rm d,*}=\begin{cases}
 \pi/2,&4b<x_*,\\
 \arcsin[x_*/(4b)],&4b\ge x_*,
 \end{cases}\qquad
 L_* =\frac{v_0s_{\rm d,*}}{\omega_{\rm r}}.
\end{equation*}
弱调制时可用 $J_1(x)\simeq x/2$，故 $|B_1|\simeq2b|\sin s_{\rm d}|$，首个最优位置是 $L_T/4$。强调制时展开 $\arcsin x\simeq x$，得
$L_*\simeq v_0x_*/(4b\omega_{\rm r})$，最优距离随场幅近似反比缩短。

达到 $x_*$ 后，第一谐波上限为 $|B_1|=J_1(x_*)\simeq0.581865$；对应 $\mathcal M_1=2|B_1|>1$。这不违反密度正性，因为 $\mathcal M_1$ 是 Fourier 调制幅度。若只保留 $1+2\Re(B_1\ee^{\ii\theta_\zeta})$，反而会出现负值，说明必须保留更高谐波。第一谐波最大、完整脉冲最窄和 Talbot 复现是三个不同条件。
\end{exbox}

到达时间抖动会改变不同次电子相对于驱动的相位。设逐发误差 $\delta t$ 为零均值高斯变量，标准差为 $\sigma_{t,\rm jit}$，且仅引起相位平移、不改变脉冲包络耦合，则
\begin{equation}
 \overline B_d=B_d\,\langle\ee^{\ii d\omega\delta t}\rangle
 =B_d\exp[-d^2\omega^2\sigma_{t,\rm jit}^2/2].
 \label{eq:electron-bunching-jitter-transfer}
\end{equation}
这就是为何高阶时间结构比总能谱更容易消失：抖动不改变 $P_\ell$，却按 $d^2$ 更快抹去高阶相干。探测器自身的高斯时间响应也给相同形式的谐波衰减，但它只模糊记录，不能据此断言到达探测器前的量子态已经退相干。

\begin{exbox}[从所需谐波反推时间抖动容差]
要求第 $d\ne0$ 阶谐波至少保留理想幅度的比例 $0<\eta_d<1$。由高斯衰减因子取对数，直接得到
\begin{equation*}
 \exp[-d^2\omega^2\sigma_{t,\rm jit}^2/2]\ge\eta_d
 \ \Longleftrightarrow\ 
 \sigma_{t,\rm jit}\le\frac{\sqrt{2\ln(1/\eta_d)}}{|d|\omega}.
\end{equation*}
因此分辨越高的谐波，允许的抖动按 $1/|d|$ 缩小。若给定同一保留阈值 $\eta$，可用谐波阶数的上限为
\begin{equation*}
 d_{\max}=\left\lfloor
 \frac{\sqrt{2\ln(1/\eta)}}{\omega\sigma_{t,\rm jit}}
 \right\rfloor,
\end{equation*}
其中向下取整表示取不超过括号内实数的最大整数。

若探测器另有独立高斯时间响应宽度 $\sigma_{t,\rm det}$，两个特征函数相乘，测得谐波为
\begin{equation*}
 B_d^{\rm obs}=B_d\exp\!\left[-\frac{d^2\omega^2}{2}
 (\sigma_{t,\rm jit}^2+\sigma_{t,\rm det}^2)\right].
\end{equation*}
因此仅凭谐波衰减，不能区分束流到达抖动与探测时间模糊。若 $B_d\ne0$，不同 $d$ 的 $-\ln|B_d^{\rm obs}/B_d|/d^2$ 应为同一常数；偏离这个关系说明单一高斯模糊模型不充分。这给出不依赖某组特定时间数值的实验检查。
\end{exbox}

\figref{fig:electron-bunching-worked} 比较相干态、相位混合以及有限抖动后的计算结果。
\begin{figure}[htbp]
 \centering
 \includegraphics[width=0.92\linewidth]{electron_bunching_worked.pdf}
 \caption{能谱到时间密度需要非对角相干。数值图取 $|\beta|=1.20$、$\SI{200}{\kilo\electronvolt}$ 电子及 $\SI{800}{\nano\metre}$ 同步光学步长。左图用相同的 PINEM 占据比较纯态、完全相位混合和有限到达抖动在 $L_*$ 的周期密度；右图给出前三阶聚束因子随漂移的变化。实空间相位与探测面到达延迟满足 $\theta_\zeta=-\omega\tau_{\rm d}$。}
 \label{fig:electron-bunching-worked}
\end{figure}

\begin{derivation}[从副对角线级数到 Bessel 闭式]
令 $x=2|\beta|$。先把\eqnrefs{eq:pinem_bunching_exact_series_dp46} 中的阶差固定为 $d$，而不截断其中的边带和：
\begin{equation*}
 B_d=\ee^{\ii d\varphi_\beta+\ii d^2s_{\rm d}}
 \sum_{\ell\in\mathbb Z}J_\ell(x)J_{\ell-d}(x)
 \ee^{-2\ii\ell ds_{\rm d}}.
\end{equation*}
对两个实 Bessel 系数分别使用互为共轭的 Fourier 表示，
\begin{equation*}
 J_\ell(x)=\frac1{2\pi}\int_0^{2\pi}\ee^{\ii x\sin u-\ii\ell u}\dd u,
 \qquad
 J_{\ell-d}(x)=\frac1{2\pi}\int_0^{2\pi}
 \ee^{-\ii x\sin v+\ii(\ell-d)v}\dd v.
\end{equation*}
此时所有 $\ell$ 依赖都集中在
$\sum_\ell\exp[-\ii\ell(u-v+2ds_{\rm d})]$。
它是 $2\pi$ 倍的周期 Dirac $\delta$，要求 $v=u+2ds_{\rm d}$ 模 $2\pi$；该等式按周期分布理解，也可先加 $\ee^{-\epsilon\ell^2}$ 收敛因子再令 $\epsilon\to0^+$。因此双角积分只剩一个。使用
\begin{equation*}
 \sin u-\sin(u+2ds_{\rm d})
 =-2\sin(ds_{\rm d})\cos(u+ds_{\rm d}),
\end{equation*}
并把剩余角平移为 $u+ds_{\rm d}$，外部的 $\ee^{\ii d^2s_{\rm d}}$ 恰与平移产生的相位抵消，得到
\begin{equation*}
 B_d=\frac{\ee^{\ii d\varphi_\beta}}{2\pi}
 \int_0^{2\pi}\dd u\,\ee^{-\ii du}
 \exp[-2\ii x\sin(ds_{\rm d})\cos u].
\end{equation*}
使用 $\exp(-\ii y\cos u)=\sum_r(-\ii)^rJ_r(y)\ee^{\ii ru}$，角积分只保留 $r=d$，便得到\eqnrefs{eq:electron-pinem-bunching-closed}。最后对抖动相位的高斯特征函数取平均，给\eqnrefs{eq:electron-bunching-jitter-transfer}。这两个步骤分别使用确定性自由传播和不同实验重复之间的统计平均，不能把后者当作额外的确定性传播相位。
\end{derivation}
